% Решения упражнений § 4 (с. 28--29).

\begin{solution}[1 (\S\,4)]
Чтобы выражение $h^{-1}[h(x)*h(y)]$ имело смысл, нужно, чтобы
$h(x)*h(y)$ лежало в области значений $h$; будем считать, что множество
$h(\R)$ замкнуто относительно $*$ (например, $h$ "--- биекция $\R\to\R$).
Строгая монотонность обеспечивает существование $h^{-1}$ на $h(\R)$
и тождество $h(h^{-1}(u))=u$. Тогда
\[
\varphi(\varphi(x,y),z)=h^{-1}\bigl[h\bigl(h^{-1}[h(x)*h(y)]\bigr)*h(z)\bigr]
=h^{-1}\bigl[(h(x)*h(y))*h(z)\bigr],
\]
\[
\varphi(x,\varphi(y,z))=h^{-1}\bigl[h(x)*h\bigl(h^{-1}[h(y)*h(z)]\bigr)\bigr]
=h^{-1}\bigl[h(x)*(h(y)*h(z))\bigr],
\]
и правые части совпадают по ассоциативности $*$. (Это <<перенос
структуры>>: $h$ "--- изоморфизм операции $\varphi$ на операцию $*$; так же
переносится и коммутативность.)

\emph{Замечание.} Все ассоциативные операции $\varphi_0$--$\varphi_3$ из
текста параграфа имеют такой вид:
\begin{itemize}[nosep]
\item $\varphi_0$: $h(x)=\operatorname{arth}(x/c)$, $*=+$ (сложение
  <<быстрот>>; $\operatorname{th}(a+b)=\frac{\operatorname{th}a+\operatorname{th}b}{1+\operatorname{th}a\operatorname{th}b}$);
\item $\varphi_1$ ($\alpha\ne0$): $h(x)=1+\alpha x$, $*=\times$, так как
  $(1+\alpha x)(1+\alpha y)=1+\alpha(x+y+\alpha xy)$;
\item $\varphi_2$ (на $x\ge0$): $h(x)=x^2+a^2$, $*=+$;
\item $\varphi_3$: $h(x)=1/x-\alpha$, $*=+$, так как
  $1/\varphi_3=1/x_1+1/x_2-\alpha$.
\end{itemize}
\end{solution}

\begin{solution}[2 (\S\,4)]
Простейший пример "--- левая проекция $x*y=x$:
$(x*y)*z=x=x*(y*z)$, но $1*2=1\ne2=2*1$.

Менее тривиальный пример: $x*y=x\,|y|$. Ассоциативность:
$(x*y)*z=x|y|\,|z|$, $x*(y*z)=x\,\bigl|y|z|\bigr|=x|y||z|$. Некоммутативность:
$1*(-1)=1$, $(-1)*1=-1$. По упр.~1 из него получается целое семейство
ассоциативных некоммутативных операций $h^{-1}[h(x)\,|h(y)|]$.
\end{solution}

\begin{solution}[3 (\S\,4)]
При $0\le t\le1$ имеем $h_t'(x)=(1-t)e^x+t>0$, так что $h_t$ строго
возрастает. Область значений: при $t=0$ это $(0,\infty)$, при $0<t\le1$
"--- все $\R$ (при $x\to-\infty$ $h_t\sim tx\to-\infty$, при $x\to+\infty$
$h_t\to+\infty$). В обоих случаях она замкнута относительно умножения,
поэтому по упр.~1 (с $*=\times$) операция
$\varphi_t(x,y)=h_t^{-1}[h_t(x)h_t(y)]$ ассоциативна, а коммутативна,
так как коммутативно умножение.

Крайние значения: $h_0=e^x$, $h_0^{-1}=\ln$,
$\varphi_0=\ln(e^xe^y)=x+y$; $h_1(x)=x$, $\varphi_1=xy$. Таким образом,
семейство $\varphi_t$ непрерывно деформирует сложение в умножение,
оставаясь в классе коммутативных ассоциативных операций.

При $0<t<1$ обратная функция выражается через функцию Ламберта $W$
($we^w=s$): из $y=(1-t)e^x+tx$, полагая $x=y/t-w$, получаем
$we^w=\frac{1-t}{t}e^{y/t}$, т.\,е.
\[
h_t^{-1}(y)=\frac yt-W\Bigl(\frac{1-t}{t}\,e^{y/t}\Bigr).
\]
\end{solution}

\begin{solution}[4 (\S\,4)]
\emph{Существование $Q(A)$.} Множество всех правильных отображений
$\mathfrak F(\mathfrak X)$ замкнуто и содержит $A$, а пересечение любого
семейства замкнутых множеств замкнуто (доказательство теоремы~1 дословно
переносится на любое число множеств). Поэтому $Q(A)$ "--- пересечение всех
замкнутых множеств, содержащих $A$.

\emph{$\M(H)$ замкнуто.} Суперпозиции суперпозиций функций из $H$ сами
являются суперпозициями функций из $H$: $\M(\M(H))=\M(H)$ (пункты
а)--г) определения применяются к уже построенным суперпозициям).

\emph{Равенство $Q(A)=\M(H)$ при $H\subset A$, $A\subset\M(H)$.}
(1) $\M(H)$ замкнуто и содержит $A$, значит, $Q(A)\subset\M(H)$.
(2) $H\subset A\subset Q(A)$ и $Q(A)$ замкнуто, значит,
$\M(H)\subset Q(A)$.

\emph{Случай $H\cap\ov A\ne\emptyset$} (система полна по отношению к $A$,
но содержит функции не из $A$). Рассуждение (1) остается в силе:
$Q(A)\subset\M(H)$; рассуждение (2) "--- нет, и включение может быть
строгим.

\emph{Примеры.} $A=C^R$ (полиномы), $H=H_0=\{x_1+x_2,\ x_1x_2,\ a\}$:
$H\subset A$, $Q(C^R)=\M(H_0)=C^R$ (множество $C^R$ само замкнуто).
$A=\{x_1+x_2,\ x_1x_2\}\cup\R$ (константы): $Q(A)=C^R$ "--- замкнутая
оболочка конечного набора операций может быть гораздо шире самого набора.
$A=C^R$, $H=H_0\cup\{\sin x\}$: $H\cap\ov A=\{\sin x\}$, и
$Q(A)=C^R\subsetneq\M(H)\ni\sin x$.
\end{solution}

\begin{solution}[5 (\S\,4)]
Первое включение "--- определение полноты: $A_i\subset\M(H_i)$ для
каждого $i$. Второе следует из монотонности операции $\M$: если
$H'\subset H''$, то всякая суперпозиция функций из $H'$ есть суперпозиция
функций из $H''$, т.\,е. $\M(H')\subset\M(H'')$; применяем это
к $H_i\subset\bigcup_jH_j$.

Включения, вообще, строгие: при $H_1=\{x_1+x_2\}$, $H_2=\{x_1x_2\}$
функция $x_1x_2+x_3$ принадлежит $\M(H_1\cup H_2)$, но не принадлежит
$\M(H_1)\cup\M(H_2)$.
\end{solution}

\begin{solution}[6 (\S\,4)]
\emph{Пример.} $\psi(x_1,x_2,x_3)=x_1-x_2+x_3$ на $\R$. Все три способа
расстановки дают одно и то же:
\begin{multline*}
\psi(\psi(x_1,x_2,x_3),x_4,x_5)=\psi(x_1,\psi(x_2,x_3,x_4),x_5)={}\\
=\psi(x_1,x_2,\psi(x_3,x_4,x_5))=x_1-x_2+x_3-x_4+x_5 .
\end{multline*}
(В любой группе аналогично $\psi=x_1x_2^{-1}x_3$; это <<груда>> Прюфера
и Бэра.)

\emph{Эта операция не имеет вида $\varphi(x_1,\varphi(x_2,x_3))$
с ассоциативной $\varphi$.} Допустим противное:
$\varphi(x,\varphi(y,z))=\varphi(\varphi(x,y),z)=x-y+z$. Пусть $R$ "---
множество значений $\varphi$. Если $u=\varphi(y,z)\in R$, то
$\varphi(x,u)=x+(z-y)$ при всех $x$; значит, $\varphi(x,u)=x+k(u)$ для
$u\in R$ (величина $k(u)=z-y$ определена однозначно). Аналогично
$\varphi(v,z)=z+m(v)$ для $v\in R$. При $u,v\in R$ отсюда
$v+k(u)=\varphi(v,u)=u+m(v)$, т.\,е. $u-k(u)=v-m(v)=c$ "--- постоянная, и
$\varphi(x,u)=x+u-c$ при $u\in R$. Тогда
$x-y+z=\varphi(x,\varphi(y,z))=x+\varphi(y,z)-c$, т.\,е.
$\varphi(y,z)=z-y+c$ при всех $y,z$; следовательно, $R=\R$ и одновременно
$\varphi(x,u)=x+u-c$ и $\varphi(x,u)=u-x+c$ при всех $x,u$, что
невозможно. Противоречие.

\emph{Обобщение на $n$-местные операции.} Операция $\psi$ от $n$
аргументов \emph{ассоциативна}, если при подстановке $\psi$ в любое из $n$
мест получается одна и та же функция $2n-1$ аргументов:
\begin{multline*}
\psi(\psi(x_1,\dots,x_n),x_{n+1},\dots,x_{2n-1})\equiv
\psi(x_1,\psi(x_2,\dots,x_{n+1}),\dots,x_{2n-1})\equiv\dots\\
\dots\equiv\psi(x_1,\dots,x_{n-1},\psi(x_n,\dots,x_{2n-1})).
\end{multline*}
Тривиальные примеры дают <<итерации>> двухместной ассоциативной операции,
$\psi=x_1*x_2*\dots*x_n$. Нетривиальные: пусть $\zeta$ "--- корень
степени $n-1$ из единицы ($\zeta^{n-1}=1$). Тогда
\[
\psi(x_1,\dots,x_n)=x_1+\zeta x_2+\zeta^2x_3+\dots+\zeta^{n-1}x_n
\]
ассоциативна: при подстановке внутреннего $\psi$ на $k$-е место его
$j$-й аргумент получает множитель $\zeta^{k-1}\zeta^{j-1}=\zeta^{(k+j-1)-1}$
"--- ровно тот, что положен месту $k+j-1$, а аргументы, сдвинувшиеся на
$n-1$ мест вправо, сохраняют множитель, так как $\zeta^{n-1}=1$.
В действительных числах можно взять $\zeta=-1$ при нечетном $n$:
$\psi=x_1-x_2+x_3-\dots+x_n$ (при $n=3$ "--- пример выше). Вообще, если
$h$ "--- биекция, то $h^{-1}[\psi(h(x_1),\dots,h(x_n))]$ снова
ассоциативна (как в упр.~1). Систематически $n$-местные ассоциативные
операции изучаются в теории $n$-арных групп (Дёрнте, 1928; Пост, 1940).
\end{solution}

\begin{solution}[7 (\S\,4)]
Функции из базиса должны принадлежать самому множеству $\M_0$, а для
проверки того, что система $H$ "--- базис, достаточно убедиться, что она
полна и что после удаления \emph{любого одного} элемента полнота
теряется. Заметим главное: при суперпозиции полиномов коэффициенты
результата лежат в кольце, порожденном (сложением и умножением)
коэффициентами исходных полиномов.

\textbf{Целые коэффициенты} ($\Z[x_1,x_2,\dots]$):
\[
H=\{x_1+x_2,\ x_1x_2,\ -1\}.
\]
Полнота: $1=(-1)(-1)$, $0=1+(-1)$, $n=1+\dots+1$, $-n=(-1)\,n$; далее
одночлены "--- произведениями, многочлены "--- суммами. Минимальность:
без $-1$ нет ни одной константы (из $x_1+x_2$ и $x_1x_2$ без констант
получаются только многочлены с неотрицательными коэффициентами без
свободного члена); без $x_1x_2$ получаются лишь функции вида
$c+\sum n_ix_i$ степени $\le1$; без $x_1+x_2$ "--- лишь одночлены
$\pm x_1^{k_1}\cdots x_s^{k_s}$ и константы $\pm1$.

\textbf{Рациональные коэффициенты} ($\Q[x_1,x_2,\dots]$):
\[
H=\{x_1+x_2,\ x_1x_2,\ -1\}\cup\{1/p:\ p\ \text{простое}\}.
\]
Полнота: кольцо, порожденное $-1$ и всеми $1/p$, есть $\Q$
($1/q$ для составного $q$ "--- произведение $1/p$). Минимальность: без
$1/p$ все коэффициенты лежат в кольце $\Z[1/q:\ q\ne p]$, где у
знаменателей нет множителя $p$, так что $1/p$ не получить; без $-1$ все
коэффициенты неотрицательны; без $x_1x_2$ или $x_1+x_2$ "--- как выше.
\emph{Конечного} базиса здесь нет: у конечной системы конечное число
коэффициентов, их знаменатели делят некоторое $N$, и все суперпозиции
имеют коэффициенты из $\Z[1/N]\ne\Q$.

\textbf{Действительные коэффициенты} ($C^R$). Система
$H_0=\{x_1+x_2,\ x_1x_2,\ a,\ \forall a\in\R\}$ полна (пример~1), но
базисом в смысле определения не является: отдельные константы избыточны
(например, $2=1+1$). Если, как в обозначении $H_0$ книги, считать семейство
всех констант одним элементом системы, то $H_0$ "--- базис: без сложения,
без умножения или без констант полноты нет (рассуждения те же, что выше).
Если же требовать неизбыточности поэлементно, то по замечанию о
коэффициентах базис вида $\{x_1+x_2,\ x_1x_2\}\cup C$ существует тогда
и только тогда, когда у кольца $\R$ есть неизбыточная система
образующих $C$ (такая, что ни одно $c\in C$ не выражается через остальные
с помощью сложения и умножения). Для $\Z$ и $\Q$ такие системы указаны
выше ($\{-1\}$ и $\{-1\}\cup\{1/p\}$); для $\R$ этот теоретико-кольцевой
вопрос выходит за рамки книги, и мы оставляем его открытым (в предисловии
автор предупреждает, что часть задач носит проблемный характер).
\end{solution}
